%-------------------------------------------------------------------------------
% cvprojects / \cvproject - the Open Source Projects section
%
% Usage: \begin{cvprojects} \cvproject{<name>}{<year>}{<url>} ... \end{cvprojects}
%
% Each project is a name and a linked URL on the left, the year on the right,
% followed by its own cvitems block and an optional \cvtechskills line.
% Pass an empty year for an entry that is just a pointer, e.g. {More Projects}{}.
%-------------------------------------------------------------------------------
% Inside \cventry the bullets get their size from \descriptionstyle, which wraps
% the description argument. There is no such wrapper here, so the environment
% applies the same switches itself - otherwise the project bullets would come
% out at the 11pt document size instead of 9pt.
\newenvironment{cvprojects}{%
	\vspace{\acvSectionContentTopSkip}
	\cvfirstentrytrue
	\begin{center}
		\fontsize{9pt}{9.96pt}\color{text}\bodyfontlight\upshape
		}{%
	\end{center}
	\cvendblock
}

\newcommand*{\projecturlstyle}[1]{{\fontsize{8pt}{1em}\bodyfontlight\color{graytext} #1}}

% The gap between the name and the link is one em of the 11pt body font. It is
% pinned to the document default family rather than inherited, because the
% surrounding environment
% runs at 9pt, and Barlow's em differs fractionally from the default family's.
\newcommand*{\projectgap}{{\normalsize\normalfont\quad}}

% The link text drops the scheme (the link itself keeps it), so the line reads
% github.com/user/repo rather than https://github.com/user/repo. gsub is used
% rather than an anchored pattern because % starts a comment inside \directlua.
\newcommand*{\stripurlscheme}[1]{%
	\directlua{
		local u = "\luaescapestring{#1}"
		u = string.gsub(u, "https://", "", 1)
		u = string.gsub(u, "http://", "", 1)
		tex.sprint(-2, u)
	}%
}

% A URL is tokenised before any macro body gets the chance to protect it, so the
% catcodes have to be switched off *before* the arguments are read. Hence the
% split: \cvproject opens the group and \cvprojectrow reads the arguments under
% it. \endgroup at the end of the row restores the normal catcodes.
\newcommand*{\cvproject}{%
	\begingroup
	\catcode`\_=12 \catcode`\~=12 \catcode`\#=12 \catcode`\&=12 \catcode`\$=12
	\cvprojectrow
}
\newcommand*{\cvprojectrow}[3]{%
	\ifcvfirstentry
		\cvtopskip{\cvfirstentrygap}%
		\global\cvfirstentryfalse
	\else
		\cvtopskip{\cvprojectgap}%
	\fi
	\setlength\tabcolsep{0pt}
	\setlength{\extrarowheight}{0pt}
	\begin{tabular*}{\textwidth}{@{\extracolsep{\fill}} L{\textwidth - 4.5cm} R{4.5cm}}
		\entrytitlestyle{#1}\projectgap\projecturlstyle{\href{#3}{\faLink\ \stripurlscheme{#3}}} & \entrydatestyle{#2} \\
	\end{tabular*}%
	\endgroup
}
